\chapter[17]{Entropy, \texorpdfstring{$\mu$}{mu}-Invariant, and Finite Time Singularities}
\label{ch:chow-iii-entropy-mu-finite-time-singularities}
\dcref{ch17}
\source{chow-et-al-ricci-flow-part-iii:page-0022}

Unless stated otherwise, \((\mathcal M^n,g)\) is closed and
\(d\mu=d\mu_g\).

\section{Perelman's entropy and its invariants}

\begin{definition}[Perelman's energy and $\lambda$-invariant]
\label{def:chow-iii-ch17-perelman-energy-lambda}
\source{chow-et-al-ricci-flow-part-iii:page-0023}
\dcref{ch17:17.1--17.4}
\group{Entropy Functionals}

For $f\in C^\infty(\mathcal M)$, put
\[
 \mathcal F(g,f)=\int_{\mathcal M}(R+|\nabla f|^2)e^{-f}\,d\mu
 =\int_{\mathcal M}(R+\Delta f)e^{-f}\,d\mu.
\]
Writing $v=e^{-f/2}$, set
\[
 \mathcal G(g,v)=\int_{\mathcal M}(Rv^2+4|\nabla v|^2)\,d\mu,
 \qquad
 \lambda(g)=\inf_{\int v^2d\mu=1}\mathcal G(g,v).
\]
The constrained minimizer is unique and smooth, and its logarithmic density
$f_0$ satisfies
\[
 2\Delta f_0-|\nabla f_0|^2+R=\lambda(g).
\]
\end{definition}

\begin{definition}[Perelman's entropy, $\mu$, and $\nu$]
\label{def:chow-iii-ch17-perelman-entropy-mu-nu}
\source{chow-et-al-ricci-flow-part-iii:page-0023,chow-et-al-ricci-flow-part-iii:page-0024}
\dcref{ch17:17.5--17.10}
\group{Entropy Functionals}

For $\tau>0$ and
$u=(4\pi\tau)^{-n/2}e^{-f}=w^2$, define
\[
 \mathcal W(g,f,\tau)=\int_{\mathcal M}
 \bigl(\tau(R+|\nabla f|^2)+f-n\bigr)u\,d\mu
\]
and equivalently
\[
 \mathcal K(g,w,\tau)=\int_{\mathcal M}
 \left[\tau(Rw^2+4|\nabla w|^2)
 -\left(\log w^2+\frac n2\log(4\pi\tau)+n\right)w^2\right]d\mu.
\]
Then
\[
 \mu(g,\tau)=\inf_{\int u\,d\mu=1}\mathcal W(g,f,\tau)
 =\inf_{\int w^2d\mu=1}\mathcal K(g,w,\tau),
 \qquad \nu(g)=\inf_{\tau>0}\mu(g,\tau).
\]
\end{definition}

\begin{theorem}[Entropy monotonicity]
\label{thm:chow-iii-ch17-entropy-monotonicity}
\source{chow-et-al-ricci-flow-part-iii:page-0024}
\uses{def:chow-iii-ch17-perelman-entropy-mu-nu}
\dcref{ch17:17.11--17.13}
\group{Entropy Functionals}


If
\[
 \partial_tg=-2\Ric,\qquad
 \partial_tf=-\Delta f+|\nabla f|^2-R+\frac n{2\tau},
 \qquad \tau'=-1,
\]
then
\[
 \frac d{dt}\mathcal W(g,f,\tau)=2\tau\int_{\mathcal M}
 \left|\Ric+\nabla^2f-\frac1{2\tau}g\right|^2u\,d\mu\ge0.
\]
\end{theorem}
\begin{proof}
\sketch For
$v=[\tau(R+2\Delta f-|\nabla f|^2)+f-n]u$, direct differentiation gives
\[
 \Box^*v=-2\tau\left|\Ric+\nabla^2f-\frac1{2\tau}g\right|^2u,
 \qquad \Box^*=-\partial_t-\Delta+R.
\]
Since $\mathcal W=\int v\,d\mu$ and $\partial_td\mu=-R\,d\mu$,
integration by parts yields $\frac d{dt}\mathcal W=-\int\Box^*v\,d\mu$.
\end{proof}

\begin{lemma}[Euler--Lagrange equation for entropy minimizers]
\label{lem:chow-iii-ch17-entropy-minimizer-equation}
\source{chow-et-al-ricci-flow-part-iii:page-0024}
\uses{def:chow-iii-ch17-perelman-entropy-mu-nu}
\dcref{ch17:eq-17.14,eq-17.15}
\group{Entropy Functionals}


A smooth constrained minimizer $f_\tau$ satisfies
\[
 \tau(2\Delta f_\tau-|\nabla f_\tau|^2+R)+f_\tau-n=\mu(g,\tau).
\]
Equivalently, for $w_\tau=(4\pi\tau)^{-n/4}e^{-f_\tau/2}$,
\[
 \tau(-4\Delta w_\tau+Rw_\tau)-w_\tau\log w_\tau^2
 -\left(\frac n2\log(4\pi\tau)+n\right)w_\tau
 =\mu(g,\tau)w_\tau.
\]
\end{lemma}
\begin{proof}
\sketch Differentiate $\mathcal K(g,w,\tau)$ along variations tangent to the unit
sphere in $L^2$.  The Lagrange multiplier is the value of $\mathcal K$ at
the minimizer, hence $\mu(g,\tau)$.  Substituting
$w=(4\pi\tau)^{-n/4}e^{-f/2}$ gives the first equation.
\end{proof}

\section{Bounds for the \texorpdfstring{$\mu$}{mu}-invariant}

\begin{lemma}[Logarithmic Sobolev inequality]
\label{lem:chow-iii-ch17-logarithmic-sobolev}
\source{chow-et-al-ricci-flow-part-iii:page-0025}
\dcref{ch17:17.1}
\group{Entropy Bounds}

For every $a>0$ there is $C(a,g)$ such that every positive $\varphi$ with
$\int\varphi^2d\mu=1$ satisfies
\[
 \int_{\mathcal M}\varphi^2\log\varphi\,d\mu
 \le a\int_{\mathcal M}|\nabla\varphi|^2\,d\mu+C(a,g),
\]
where, for $n\ge3$,
\[
 C(a,g)=a\Vol(g)^{-2/n}+\frac{n^2}{4ae^2C_s(\mathcal M,g)}.
\]
\end{lemma}

\begin{lemma}[Sobolev constants under scaling the metric]
\label{lem:chow-iii-ch17-sobolev-scaling}
\source{chow-et-al-ricci-flow-part-iii:page-0025,chow-et-al-ricci-flow-part-iii:page-0026}
\dcref{ch17:17.2}
\group{Entropy Bounds}

For $\lambda>0$,
$C_s(\mathcal M,\lambda^2g)=C_s(\mathcal M,g)$; if $\lambda\ge1$, then
$C(a,\lambda^2g)\le C(a,g)$.
\end{lemma}
\begin{proof}
\sketch Put $\widetilde g=\lambda^2g$ and
$\widetilde\varphi=\lambda^{-n/2}\varphi$.  The normalized $L^2$ measure is
unchanged, while both terms in the defining Sobolev inequality acquire the
factor $\lambda^{-2}$, proving the first identity.  The displayed formula
for $C(a,g)$, volume scaling, and the first identity prove the second.
\end{proof}

\begin{lemma}[Upper bound for $\mu$ in terms of $\lambda$, volume, $\tau$, and $n$]
\label{lem:chow-iii-ch17-mu-upper-lambda-volume}
\source{chow-et-al-ricci-flow-part-iii:page-0026}
\uses{def:chow-iii-ch17-perelman-energy-lambda, def:chow-iii-ch17-perelman-entropy-mu-nu}
\dcref{ch17:17.3}
\group{Entropy Bounds}


For every $\tau>0$,
\[
 \mu(g,\tau)\le\tau\lambda(g)+\frac1e\Vol(g)
 -\frac n2\log(4\pi\tau)-n.
\]
\end{lemma}
\begin{proof}
\sketch For $\int w^2d\mu=1$, the inequality
$w^2\log w^2\ge-e^{-1}$ gives
\[
 \mathcal K(g,w,\tau)\le \tau\mathcal G(g,w)+e^{-1}\Vol(g)
 -\frac n2\log(4\pi\tau)-n.
\]
Evaluate this at the constrained minimizer of $\mathcal G$.
\end{proof}

\begin{corollary}[Upper bound for $\mu$ when $\lambda\le0$]
\label{cor:chow-iii-ch17-mu-upper-nonpositive-lambda}
\source{chow-et-al-ricci-flow-part-iii:page-0026,chow-et-al-ricci-flow-part-iii:page-0027}
\uses{lem:chow-iii-ch17-mu-upper-lambda-volume}
\dcref{ch17:17.4}
\group{Entropy Bounds}


If $\lambda(g)\le0$, then
\[
 \mu(g,\tau)\le\log\Vol(g)-\frac n2\log(4\pi\tau)-n+\frac1e.
\]
\end{corollary}
\begin{proof}
\sketch Use $\mu(g,\tau)=\mu(cg,c\tau)$ in the preceding lemma with
$c=\Vol(g)^{-2/n}$, and use $\lambda(cg)=c^{-1}\lambda(g)\le0$.
\end{proof}

\begin{lemma}[Lower bound for $\mu$]
\label{lem:chow-iii-ch17-mu-lower-log-sobolev}
\source{chow-et-al-ricci-flow-part-iii:page-0027}
\uses{lem:chow-iii-ch17-logarithmic-sobolev}
\dcref{ch17:17.5}
\group{Entropy Bounds}


For every $\tau>0$,
\[
 \mu(g,\tau)\ge \tau R_{\min}(g)-2C(2\tau,g)
 -\frac n2\log(4\pi\tau)-n.
\]
\end{lemma}
\begin{proof}
\sketch Apply Lemma~\ref{lem:chow-iii-ch17-logarithmic-sobolev} with $a=2\tau$
to a normalized $w$.  In the formula for $\mathcal K$, its gradient term
cancels the resulting entropy estimate, and
$\int Rw^2d\mu\ge R_{\min}(g)$.  Take the infimum over $w$.
\end{proof}

\begin{lemma}[Lower bound for the $\mu$-invariant]
\label{lem:chow-iii-ch17-mu-lower-lambda}
\source{chow-et-al-ricci-flow-part-iii:page-0028,chow-et-al-ricci-flow-part-iii:page-0029}
\dcref{ch17:17.6}
\group{Entropy Bounds}

If $\tau\ge n/8$, then
\[
 \mu(g,\tau)\ge\left(\tau-\frac n8\right)\lambda(g)
 -\frac n2\log(4\pi\tau)-n-n\log C_s(g)+\frac n8R_{\min}(g).
\]
\end{lemma}
\begin{proof}
\sketch For $\int w^2d\mu=1$, Jensen's inequality and the $L^2$ Sobolev inequality
give
\[
 -\int w^2\log w^2d\mu\ge-n\log C_s(g)
 -\frac n2\int|\nabla w|^2d\mu.
\]
Moreover
$4\int|\nabla w|^2d\mu\le\mathcal G(g,w)-R_{\min}(g)$.
Insert both estimates into $\mathcal K$ and use
$\mathcal G(g,w)\ge\lambda(g)$; the coefficient $\tau-n/8$ is
nonnegative.
\end{proof}

\begin{remark}[Large- and small-scale behavior of the lower bound]
\label{rem:chow-iii-ch17-mu-lower-bound-scales}
\source{chow-et-al-ricci-flow-part-iii:page-0028}
\uses{lem:chow-iii-ch17-mu-lower-lambda}
\dcref{ch17:17.7}
\group{Entropy Bounds}


The preceding estimate cannot extend to all sufficiently small $\tau$, since
its right side tends to $+\infty$ as $\tau\downarrow0$.  If
$\lambda(g)>0$, it does imply $\mu(g,\tau)\to+\infty$ as
$\tau\to\infty$.
\end{remark}

\begin{lemma}[Lower bound for the volume of a solution when $\lambda\le0$]
\label{lem:chow-iii-ch17-volume-lower-nonpositive-lambda}
\source{chow-et-al-ricci-flow-part-iii:page-0029}
\uses{cor:chow-iii-ch17-mu-upper-nonpositive-lambda, lem:chow-iii-ch17-mu-lower-lambda, thm:chow-iii-ch17-entropy-monotonicity}
\dcref{ch17:17.8}
\group{Finite Time Singularities}


If $g(t)$, $t\in[0,T)$, is a closed Ricci flow and
$\lambda(g(t))\le0$ throughout, then constants $c_1,c_2>0$ depending only
on $g(0)$ satisfy
\[
 \Vol(g(t))\ge c_1e^{-c_2t}.
\]
\end{lemma}
\begin{proof}
\sketch Entropy monotonicity and the lower bound at scale $n/8+t$ give
\[
 \mu(g(t),n/8)\ge t\lambda(g(0))
 -\frac n2\log(4\pi(n/8+t))+C(g(0)).
\]
The upper bound at scale $n/8$ bounds this by
$\log\Vol(g(t))+C(n)$.  Exponentiation gives a positive multiple of
$(n/8+t)^{-n/2}e^{t\lambda(g(0))}$, which on $[0,T)$ is bounded below by
$c_1e^{-c_2t}$ after enlarging $c_2$.
\end{proof}

\section{Compact finite-time singularity models}

\begin{remark}[Singularity models]
\label{rem:chow-iii-ch17-singularity-models}
\source{chow-et-al-ricci-flow-part-iii:page-0030}
\dcref{ch17:17.9}
\group{Finite Time Singularities}

A finite-time singularity model is a complete ancient pointed limit of
parabolic rescalings of a finite-time singular Ricci flow on a closed
manifold.
\end{remark}

\begin{lemma}[Finite-time singularity models are noncompact when $\lambda\le0$]
\label{lem:chow-iii-ch17-model-noncompact-nonpositive-lambda}
\source{chow-et-al-ricci-flow-part-iii:page-0030}
\uses{lem:chow-iii-ch17-volume-lower-nonpositive-lambda, rem:chow-iii-ch17-singularity-models}
\dcref{ch17:17.10}
\group{Finite Time Singularities}


If a singular closed Ricci flow has $\lambda(g(t))\le0$ throughout, every
finite-time singularity model is noncompact.
\end{lemma}
\begin{proof}
\sketch For a blow-up factor $K_i\to\infty$, the rescaled time-zero volume is
$K_i^{n/2}\Vol(g(t_i))$.  The preceding volume lower bound makes this tend
to infinity.  Smooth convergence to a compact model would instead make it
converge to that model's finite volume, a contradiction.
\end{proof}

\begin{remark}[Noncompact singularity models have infinite volume]
\label{rem:chow-iii-ch17-noncompact-model-infinite-volume}
\source{chow-et-al-ricci-flow-part-iii:page-0030}
\dcref{ch17:17.11}
\group{Finite Time Singularities}

A bounded-curvature noncompact finite-time singularity model has infinite
volume at every time.
\end{remark}
\begin{proof}
\sketch No local collapsing gives a uniform positive lower bound for the volume of
every unit ball.  A maximal disjoint family of fixed smaller balls is
infinite on a complete noncompact manifold, so their volumes sum to infinity.
\end{proof}

\begin{corollary}[A compact singularity model implies $\lambda>0$]
\label{cor:chow-iii-ch17-compact-model-positive-lambda}
\source{chow-et-al-ricci-flow-part-iii:page-0030}
\uses{lem:chow-iii-ch17-model-noncompact-nonpositive-lambda}
\dcref{ch17:17.12}
\group{Finite Time Singularities}


If a finite-time singular closed Ricci flow has a compact singularity model,
then $\lambda(g(t_0))>0$ for some $t_0<T$, and hence
$\lambda(g(t))>0$ for every $t\in[t_0,T)$.
\end{corollary}
\begin{proof}
\sketch The contrapositive of the preceding lemma gives positivity at some time.
Monotonicity of $\lambda$ preserves strict positivity thereafter.
\end{proof}

\begin{theorem}[Compact finite-time singularity models are shrinkers]
\label{thm:chow-iii-ch17-compact-model-shrinker}
\source{chow-et-al-ricci-flow-part-iii:page-0031,chow-et-al-ricci-flow-part-iii:page-0032}
\uses{cor:chow-iii-ch17-compact-model-positive-lambda, lem:chow-iii-ch17-nu-continuity-positive-lambda}
\dcref{ch17:17.13}
\group{Finite Time Singularities}


Every finite-time singularity model whose underlying manifold is closed is a
shrinking gradient Ricci soliton.
\end{theorem}
\begin{proof}
\sketch Let $g_i(t)=Q_i g(t_i+Q_i^{-1}t)$ converge smoothly to $g_\infty(t)$.
After translating time, $\lambda(g(t))>0$.  Monotonicity and negativity of
$\nu$ give a limit $\nu_T=\lim_{t\uparrow T}\nu(g(t))$.
Smooth convergence and scaling imply $\lambda(g_\infty(t))\ge0$.  Equality
at one time would, by the equality case of $\lambda$-monotonicity, make the
compact ancient model steady and hence Ricci flat.  No local collapsing at
all scales would then force its compact volume to exceed $\kappa r^n$ for
every $r$, an impossibility.  Thus $\lambda(g_\infty(t))>0$.
Lemma~\ref{lem:chow-iii-ch17-nu-continuity-positive-lambda} and scaling now
give $\nu(g_\infty(t))=\nu_T$ for all $t$.  Equality in the
$\nu$-monotonicity formula therefore gives the shrinking soliton equation.
\end{proof}

\begin{lemma}[Continuity of $\nu$ at positive $\lambda$]
\label{lem:chow-iii-ch17-nu-continuity-positive-lambda}
\source{chow-et-al-ricci-flow-part-iii:page-0032,chow-et-al-ricci-flow-part-iii:page-0033}
\uses{lem:chow-iii-ch17-mu-continuity, lem:chow-iii-ch17-mu-lower-lambda, prop:chow-iii-ch17-mu-small-scale-limit}
\dcref{ch17:17.14}
\group{Finite Time Singularities}


If $g_i\to g_\infty$ smoothly on a closed manifold and
$\lambda(g_\infty)>0$, then $\nu(g_i)\to\nu(g_\infty)$.
\end{lemma}
\begin{proof}
\sketch At one fixed scale, small-scale negativity and continuity of $\mu$ bound all
$\nu(g_i)$ above by a common negative constant.  The large-scale lower bound
for $\mu$, together with $\lambda(g_i)\to\lambda(g_\infty)>0$, confines all
minimizing scales to a common bounded interval away from infinity.  The
small-scale limit of $\mu$, uniformly for the smoothly converging metrics,
confines them away from zero.  Uniform continuity of $\mu(g_i,\tau)$ on the
remaining compact interval lets one pass the infimum to the limit.
\end{proof}

\begin{lemma}[Continuous dependence of $\mu(g,\tau)$ on $g$]
\label{lem:chow-iii-ch17-mu-continuity}
\source{chow-et-al-ricci-flow-part-iii:page-0033,chow-et-al-ricci-flow-part-iii:page-0034}
\dcref{ch17:17.15}
\group{Entropy Bounds}

For $n\ge2$, $C<\infty$, and $\varepsilon>0$, there is $\delta>0$ with
the following property.  If $R_{\rm avg}(\widetilde g),\Vol(\widetilde g)le C$,
$|R_g-R_{\widetilde g}|\le\delta$, and
$|g-\widetilde g|_{\widetilde g}\le\delta$, then
\[
 \mu(g,\tau)-\mu(\widetilde g,\tau)\le\varepsilon
 \quad\text{for }\tau\in[C^{-1},C].
\]
\end{lemma}
\begin{proof}
\sketch Let $\widetilde w$ minimize $\mathcal K(\widetilde g,\cdot,\tau)$ and
renormalize it by a constant $c\to1$ in $L^2(g)$.  Comparing
$\mathcal K(g,c\widetilde w,\tau)$ with
$\mathcal K(\widetilde g,\widetilde w,\tau)$ reduces the difference to the
metric, scalar-curvature, volume-form, gradient, and logarithmic terms.
The first three are uniformly small by the hypotheses.  The logarithmic
Sobolev inequality and the elementary upper bound for $\mu$ uniformly bound
the last two on $[C^{-1},C]$.  Choosing $\delta$ sufficiently small makes
their total difference at most $\varepsilon$.
\end{proof}

\begin{corollary}[Singularity models on closed three-manifolds are round]
\label{cor:chow-iii-ch17-closed-three-model-round}
\source{chow-et-al-ricci-flow-part-iii:page-0034}
\uses{thm:chow-iii-ch17-compact-model-shrinker}
\dcref{ch17:17.16}
\group{Finite Time Singularities}


A finite-time singularity model on a closed three-manifold is a shrinking
spherical space form.
\end{corollary}
\begin{proof}
\sketch The preceding theorem makes the model a compact shrinking gradient Ricci
soliton.  In dimension three every such soliton has constant positive
sectional curvature, and hence is a spherical space form.
\end{proof}

\begin{remark}[Type of a compact singularity model]
\label{rem:chow-iii-ch17-compact-model-type-i-problem}
\source{chow-et-al-ricci-flow-part-iii:page-0035}
\dcref{ch17:17.17}
\group{Finite Time Singularities}

The source asks whether every finite-time singular Ricci flow that forms a
singularity model on a closed manifold must be Type I.
\end{remark}

\begin{remark}[Compact factors of singularity models are shrinkers]
\label{rem:chow-iii-ch17-compact-factor-shrinker}
\source{chow-et-al-ricci-flow-part-iii:page-0036}
\uses{thm:chow-iii-ch17-compact-model-shrinker}
\dcref{ch17:17.18}
\group{Finite Time Singularities}


If the universal cover of a finite-time singularity model splits as
$(\mathcal N^m,h(t))\times\mathbb R^{n-m}$ with $\mathcal N$ compact, then
$(\mathcal N,h(t))$ is a shrinking gradient Ricci soliton.
\end{remark}

\section{Small-scale behavior of \texorpdfstring{$\mu$}{mu}}

\begin{lemma}[$\mu(g,\tau)$ is negative for $\tau$ small]
\label{lem:chow-iii-ch17-mu-negative-small-scale}
\source{chow-et-al-ricci-flow-part-iii:page-0036,chow-et-al-ricci-flow-part-iii:page-0037}
\uses{thm:chow-iii-ch17-entropy-monotonicity}
\dcref{ch17:17.19}
\group{Small Scale Entropy}


There is $\bar\tau>0$ such that $\mu(g,\tau)<0$ for every
$0<\tau<\bar\tau$.
\end{lemma}
\begin{proof}
\sketch Run Ricci flow from $g$ to time $\bar\tau$ and use the adjoint heat kernel
based at $(x_0,\bar\tau)$.  Its entropy tends to zero at the pole, so entropy
monotonicity gives $\mu(g,\bar\tau)\le0$.  Equality would force
$\Ric+\nabla^2f-g/(2\tau)=0$ throughout the interval.  Smoothness at the
terminal time would then force the shrinking soliton to be flat, impossible
on a closed manifold.  Thus the inequality is strict; repeating the argument
with any smaller terminal time proves the claim.
\end{proof}

\begin{proposition}[$\mu(g,\tau)\to0$ as $\tau\to0$]
\label{prop:chow-iii-ch17-mu-small-scale-limit}
\source{chow-et-al-ricci-flow-part-iii:page-0037,chow-et-al-ricci-flow-part-iii:page-0038,chow-et-al-ricci-flow-part-iii:page-0039,chow-et-al-ricci-flow-part-iii:page-0040,chow-et-al-ricci-flow-part-iii:page-0041,chow-et-al-ricci-flow-part-iii:page-0042}
\uses{lem:chow-iii-ch17-mu-negative-small-scale, prop:chow-iii-ch17-smooth-entropy-minimizer, lem:chow-iii-ch17-weak-strong-maximum-principle, lem:chow-iii-ch17-minimizer-maximum-lower-bound}
\dcref{ch17:17.20}
\group{Small Scale Entropy}


For every closed Riemannian manifold,
\[
 \lim_{\tau\downarrow0}\mu(g,\tau)=0.
\]
\end{proposition}
\begin{proof}
\sketch Negativity gives $\limsup\mu\le0$.  If the lower limit were negative, choose
$\tau_i\downarrow0$ with $\mu(g,\tau_i)\le-\varepsilon$ and rescale
$g_i=(2\tau_i)^{-1}g$.  Smooth entropy minimizers $w_i$ at scale $1/2$
satisfy a uniform elliptic equation and uniform $W^{1,2}$ and local
$W^{2,p}$ estimates.  In pointed exponential coordinates based at maxima of
$w_i$, a subsequence converges in $C^{1,\alpha}_{\rm loc}$ to
$0\ne w_\infty\in W^{1,2}(\mathbb R^n)$; nontriviality follows from the
maximum lower bound, and the weak strong maximum principle makes
$w_\infty>0$.  Elliptic regularity makes it smooth and gives
\[
 2\Delta w_\infty=-\left(\mu_\infty+\frac n2\log(2\pi)+n
 +2\log w_\infty\right)w_\infty,
 \qquad \mu_\infty\le-\varepsilon.
\]
Apply the Euclidean logarithmic Sobolev inequality to
$\eta_Rw_\infty$, where $\eta_R$ is a radial cutoff.  Testing the equation
against $\eta_R^2w_\infty$ shows that its nonnegative logarithmic Sobolev
functional equals
\[
 \frac{\mu_\infty}{2}\int(\eta_Rw_\infty)^2
 +\int(|\nabla\eta_R|^2-\eta_R^2\log\eta_R)w_\infty^2.
\]
The second term tends to zero as $R\to\infty$, while the first tends to a
strictly negative number, a contradiction.
\end{proof}

\begin{remark}[Small-scale entropy asymptotics]
\label{rem:chow-iii-ch17-mu-small-scale-asymptotics-problem}
\source{chow-et-al-ricci-flow-part-iii:page-0043}
\dcref{ch17:17.21}
\group{Small Scale Entropy}

The source asks for the more precise asymptotic expansion of $\mu(g,\tau)$
as $\tau\downarrow0$.
\end{remark}

\begin{lemma}[For small $\tau$ the minimizer is nonconstant and may not be unique]
\label{lem:chow-iii-ch17-small-minimizer-nonconstant-nonunique}
\source{chow-et-al-ricci-flow-part-iii:page-0043}
\uses{lem:chow-iii-ch17-entropy-minimizer-equation, prop:chow-iii-ch17-mu-small-scale-limit}
\dcref{ch17:17.22}
\group{Small Scale Entropy}


For all sufficiently small $\tau$, every minimizer of
$\mathcal W(g,\cdot,\tau)$ is nonconstant.  If the isometry group of
$(\mathcal M,g)$ acts transitively, then every such minimizer is nonunique.
\end{lemma}
\begin{proof}
\sketch If constant minimizers existed along $\tau_i\downarrow0$, the
Euler--Lagrange equation would make $R$ constant and normalization would give
\[
 \mu(g,\tau_i)=\tau_iR-\frac n2\log(4\pi\tau_i)+\log\Vol(g)-n\to+\infty,
\]
contrary to Proposition~\ref{prop:chow-iii-ch17-mu-small-scale-limit}.
For the second claim, uniqueness would make the minimizer invariant under
every isometry.  Transitivity would then make it constant, contradicting the
first claim.
\end{proof}

\begin{remark}[Minimizers on homogeneous spaces]
\label{rem:chow-iii-ch17-homogeneous-minimizer-problem}
\source{chow-et-al-ricci-flow-part-iii:page-0044}
\dcref{ch17:17.23}
\group{Small Scale Entropy}

The source asks for a description of entropy minimizers on round spheres and
on other homogeneous manifolds, including their symmetries and their change
across the Einstein scale.
\end{remark}

\section{Existence and regularity of minimizers}

\begin{proposition}[Existence of a smooth minimizer for $\mathcal W$]
\label{prop:chow-iii-ch17-smooth-entropy-minimizer}
\source{chow-et-al-ricci-flow-part-iii:page-0044,chow-et-al-ricci-flow-part-iii:page-0045,chow-et-al-ricci-flow-part-iii:page-0046}
\uses{lem:chow-iii-ch17-mu-lower-log-sobolev, lem:chow-iii-ch17-log-entropy-continuity, lem:chow-iii-ch17-weak-strong-maximum-principle}
\dcref{ch17:17.24}
\group{Entropy Minimizers}


For every metric on a closed manifold and every $\tau>0$, there exists a
smooth constrained minimizer $f_\tau$ of $\mathcal W(g,\cdot,\tau)$, and it
satisfies the entropy Euler--Lagrange equation.
\end{proposition}
\begin{proof}
\sketch By scaling take $\tau=1$.  A normalized nonnegative minimizing sequence for
$\mathcal K$ is bounded in $W^{1,2}$ by the lower entropy estimate.  Weak
compactness and Rellich compactness give a nonnegative normalized limit
$w_\infty$.  Lower semicontinuity of the Dirichlet energy and continuity of
$\int w^2\log w^2$ show that it attains the infimum.  Its constrained first
variation is the weak Euler--Lagrange equation.  Elliptic $L^p$ estimates and
Sobolev embedding give $w_\infty\in C^{1,\alpha}$.  The weak strong maximum
principle makes it positive, after which Schauder bootstrapping makes it
smooth.  Finally
$f_1=-\frac n2\log(4\pi)-2\log w_\infty$ is the desired minimizer; scaling
restores arbitrary $\tau$.
\end{proof}

\begin{lemma}[Continuity of logarithmic entropy]
\label{lem:chow-iii-ch17-log-entropy-continuity}
\source{chow-et-al-ricci-flow-part-iii:page-0047}
\dcref{ch17:17.25}
\group{Entropy Minimizers}

For every $\delta>0$, the map
\[
 w\longmapsto\int_{\mathcal M}w^2\log w^2\,d\mu
\]
is continuous in $L^{2(1+\delta)}$, and hence in $W^{1,2}$ on a closed
manifold.
\end{lemma}
\begin{proof}
\sketch The mean value theorem writes the pointwise difference as
$(|w_2|-|w_1|)2a(1+\log a^2)$ with $a$ between $|w_1|$ and $|w_2|$.
Cauchy--Schwarz applies because
$|a\log a|\le\max\{e^{-1},(\delta e)^{-1}a^{1+\delta}\}$.
Thus the difference of the integrals tends to zero with the
$L^{2(1+\delta)}$ distance.  Sobolev embedding supplies the last assertion.
\end{proof}

\begin{lemma}[Strong maximum principle for weak solutions]
\label{lem:chow-iii-ch17-weak-strong-maximum-principle}
\source{chow-et-al-ricci-flow-part-iii:page-0047,chow-et-al-ricci-flow-part-iii:page-0048,chow-et-al-ricci-flow-part-iii:page-0049,chow-et-al-ricci-flow-part-iii:page-0050,chow-et-al-ricci-flow-part-iii:page-0051}
\dcref{ch17:17.26}
\group{Entropy Minimizers}

Let $(\mathcal M,g)$ be complete and let $w\ge0$ be a $C^{1,\alpha}$ weak
solution of the entropy minimizer equation at scale one.  If $w(p)=0$, then
$w$ vanishes in a neighborhood of $p$.
\end{lemma}
\begin{proof}
\sketch Let $F(r)$ be the average of $w$ over $S(p,r)$ and $G(r)$ the average of its
radial derivative.  Smooth polar coordinates give
$F'(r)\le G(r)+CrF(r)$.  Testing the weak equation against radial functions,
and applying Jensen's inequality to $x\log x$, gives
\[
 (G(r)A(r))'\le\left(CF(r)-\frac12F(r)\log F(r)\right)A(r),
\]
where $A(r)=\Vol_{n-1}S(p,r)$.  Since $A(r)$ is comparable to $r^{n-1}$
near zero and $F(r)\to0$, integration first yields $F(r)\le r$.  If
$F(r)\le r^k$, the same integral inequality yields, on one fixed smaller
interval, $F(r)\le r^{k+1/2}$.  Iteration gives $F(r)\le r^\ell$ for every
$\ell$, hence $F=0$ there.  Since $w\ge0$, vanishing of its spherical
averages implies that $w$ vanishes on the corresponding ball.
\end{proof}

\begin{remark}[Scaling of the weak maximum principle]
\label{rem:chow-iii-ch17-weak-maximum-principle-scaling}
\source{chow-et-al-ricci-flow-part-iii:page-0048}
\uses{lem:chow-iii-ch17-weak-strong-maximum-principle}
\dcref{ch17:17.27}
\group{Entropy Minimizers}


By metric scaling, the preceding weak maximum principle holds for the
entropy minimizer equation at every scale $\tau>0$.
\end{remark}

\begin{lemma}[Lower bound for the maximum value of a minimizer]
\label{lem:chow-iii-ch17-minimizer-maximum-lower-bound}
\source{chow-et-al-ricci-flow-part-iii:page-0052}
\uses{lem:chow-iii-ch17-entropy-minimizer-equation}
\dcref{ch17:17.28}
\group{Entropy Minimizers}


For an entropy minimizer $w_\tau$ on a closed manifold,
\[
 \max_{\mathcal M}w_\tau\ge
 \exp\left(\frac\tau2R_{\min}(g)-\frac n4\log(4\pi\tau)
 -\frac n2-\frac12\mu(g,\tau)\right).
\]
\end{lemma}
\begin{proof}
\sketch At a maximum point of $w_\tau$, one has $\Delta w_\tau\le0$.  Substitute
this and $R\ge R_{\min}$ into the Euler--Lagrange equation and divide by the
positive maximum value; solving the resulting inequality for
$\log w_\tau$ gives the formula.
\end{proof}

\section{One- and two-loop variation formulas}

\begin{definition}[Linear trace Harnack quadratic]
\label{def:chow-iii-ch17-linear-trace-harnack}
\source{chow-et-al-ricci-flow-part-iii:page-0052,chow-et-al-ricci-flow-part-iii:page-0053}
\dcref{ch17:17.89--17.93}
\group{Variation Formulas}

For a symmetric two-tensor $v$ and vector field $X$, define
\[
 L(v,X)=\operatorname{div}(\operatorname{div}v)+\langle v,\Ric\rangle
 -2\langle\operatorname{div}v,X\rangle+v(X,X).
\]
For $X=\nabla f$ this is
\[
 L(v,\nabla f)=(\operatorname{div}-\nabla f)^2v
 +\langle\Ric+\nabla^2f,v\rangle.
\]
\end{definition}

\begin{lemma}[Variation of Perelman's modified scalar curvature]
\label{lem:chow-iii-ch17-modified-scalar-variation}
\source{chow-et-al-ricci-flow-part-iii:page-0053}
\uses{def:chow-iii-ch17-linear-trace-harnack}
\dcref{ch17:17.29}
\group{Variation Formulas}


If $\partial_sg=v$ and $\partial_sf=\operatorname{tr}_gv$, so that
$e^{-f}d\mu$ is fixed, then
\[
 \partial_s(R+2\Delta f-|\nabla f|^2)
 =L(v,\nabla f)-2\langle v,\Ric+\nabla^2f\rangle.
\]
\end{lemma}
\begin{proof}
\sketch Use the standard variations of scalar curvature, the Laplacian, and squared
gradient norm.  The double-divergence and Ricci contractions assemble into
$L(v,\nabla f)$; the remaining Hessian and Ricci terms are exactly the last
inner product.
\end{proof}

\begin{lemma}[Perelman's first variation formula for $\mathcal F$]
\label{lem:chow-iii-ch17-perelman-energy-first-variation}
\source{chow-et-al-ricci-flow-part-iii:page-0053}
\uses{lem:chow-iii-ch17-modified-scalar-variation}
\dcref{ch17:17.30}
\group{Variation Formulas}


Under the same variation,
\[
 \partial_s\mathcal F(g,f)=-\int_{\mathcal M}
 \langle v,\Ric+\nabla^2f\rangle e^{-f}d\mu
 =-\int_{\mathcal M}L(v,\nabla f)e^{-f}d\mu.
\]
\end{lemma}
\begin{proof}
\sketch Differentiate the modified-scalar formula for $\mathcal F$ and use that
$e^{-f}d\mu$ is fixed.  Weighted integration by parts makes the integral of
$L(v,\nabla f)$ equal to the integral of
$\langle v,\Ric+\nabla^2f\rangle$.
\end{proof}

\begin{definition}[Matrix Harnack tensor and curvature square]
\label{def:chow-iii-ch17-matrix-harnack-curvature-square}
\source{chow-et-al-ricci-flow-part-iii:page-0053,chow-et-al-ricci-flow-part-iii:page-0054}
\dcref{ch17:17.97--17.103}
\group{Variation Formulas}

Let
\[
 P(X,Y,Z)=(\nabla_X\Ric)(Y,Z)-(\nabla_Y\Ric)(X,Z)
\]
and define
\[
 H_{\nabla f}(X,Y)=
 (\Delta_L\Ric-\tfrac12\nabla^2R+\Ric^2)(X,Y)
 +P(X,\nabla f,Y)+P(Y,\nabla f,X)
 +\Rm(\nabla f,X,Y,\nabla f).
\]
Also set $\alpha_{ij}=R_{ik\ell p}R_j{}^{k\ell p}$.  Then
\[
 \operatorname{div}\left(\alpha-\frac14|\Rm|^2g\right)
 =\Rm(\,cdot\,,e_i,e_j,e_k)P^*(e_i,e_j,e_k).
\]
\end{definition}

\begin{theorem}[Coupled one- and two-loop variation identity]
\label{thm:chow-iii-ch17-coupled-loop-variation}
\source{chow-et-al-ricci-flow-part-iii:page-0054,chow-et-al-ricci-flow-part-iii:page-0055,chow-et-al-ricci-flow-part-iii:page-0056}
\uses{lem:chow-iii-ch17-perelman-energy-first-variation, def:chow-iii-ch17-matrix-harnack-curvature-square}
\dcref{ch17:17.104--17.112}
\group{Variation Formulas}


Define
\[
 \beta^{(1)}=-2(\Ric+\nabla^2f),\quad \gamma^{(1)}=-R-\Delta f,
 \quad \mathcal F_1=\int(R+|\nabla f|^2)e^{-f}d\mu,
\]
\[
 \beta^{(2)}=-\alpha,\quad \gamma^{(2)}=-\frac12|\Rm|^2,
 \quad \mathcal F_2=\frac14\int|\Rm|^2e^{-f}d\mu.
\]
If $\delta_{(\beta,\gamma)}$ denotes variation in the direction
$(\partial_sg,\partial_sf)=(\beta,\gamma)$, then
\[
 -\delta_{(\beta^{(1)},\gamma^{(1)})}\mathcal F_2
 +\delta_{(\beta^{(2)},\gamma^{(2)})}\mathcal F_1
 =\int_{\mathcal M}|P^*-\iota_{\nabla f}\Rm|^2e^{-f}d\mu\ge0.
\]
\end{theorem}
\begin{proof}
\sketch The curvature evolution and the contracted Bianchi identity for $\alpha$
give
\[
 \frac14\Box^*(|\Rm|^2e^{-f})
 =e^{-f}\bigl(-L(\alpha,\nabla f)
 +|P^*-\iota_{\nabla f}\Rm|^2\bigr).
\]
On the other hand, the modified-scalar variation formula in the direction
$(-\alpha,-|\Rm|^2/2)$ gives the same weighted
$L(\alpha,\nabla f)$ term with the opposite sign.  Integrating on the closed
manifold cancels the divergence and $L$ terms, leaving precisely the stated
square.
\end{proof}

\begin{remark}[Zeroth-loop variation identity]
\label{rem:chow-iii-ch17-zeroth-loop-variation}
\source{chow-et-al-ricci-flow-part-iii:page-0056}
\dcref{ch17:17.31}
\group{Variation Formulas}

If $\beta^{(0)}=g$, $\gamma^{(0)}=n/2$, and
$\mathcal F_0=-\int f e^{-f}d\mu$, then
\[
 \delta_{(\beta^{(0)},\gamma^{(0)})}\mathcal F_1
 +\delta_{(\beta^{(1)},\gamma^{(1)})}\mathcal F_0=0.
\]
\end{remark}

\begin{remark}[Guiding principles for Ricci flow]
\label{rem:chow-iii-ch17-guiding-principles-problem}
\dcref{ch17:17.33}
\group{Further Problems}
\source{chow-et-al-ricci-flow-part-iii:page-0057}
The source asks for new guiding principles beyond heat-equation analogies,
maximum principles, monotonicity, solitons, space-time quantities, and point
picking, with particular emphasis on localization, weak continuation through
singularities, higher dimensions, and special geometric structures.
\end{remark}

\begin{remark}[Higher-loop monotonicity problem]
\label{rem:chow-iii-ch17-higher-loop-problem}
\dcref{ch17:17.32}
\group{Variation Formulas}
\source{chow-et-al-ricci-flow-part-iii:page-0056}
The source asks whether the preceding identities extend to an infinite
sequence of tensors $\beta^{(k)}$, functions $\gamma^{(k)}$, and functionals
$\mathcal F_k$ producing a monotonicity formula extending Perelman's entropy.
\end{remark}
